\begin{abcliste}
\abc An electron accelerated through $U$ gets the energy $eU$; a photon can take all of it:
\begin{align}
 \ssc{E}{max} &= e U = \result{\EmP{3}{2}}
\end{align}

\abc
\begin{align}
 \ssc{\lambda}{min} &= \frac{h c}{\ssc{E}{max}} = \lmF \\
   &= \frac{\nch\times\ncc}{\nce\times\U} \\
   &= \lm \approx \result{\lmP{-12}{2}}
\end{align}
about a tenth of the size of an atom.

\abc At \UbO\ the spectrum reaches up to \SI{70}{\kilo\electronvolt} ($\ssc{\lambda}{min} = \lmbP{-12}{2}$): the tube emits $\result{\text{more energetic photons}}$, which are absorbed less. (More photons of the same energies alone would not get through more tissue.)
\end{abcliste}
